\documentclass[aspectratio=169,11pt]{beamer}
\usepackage[T1]{fontenc}
\usepackage{lmodern}
\usepackage{booktabs}
\usepackage{amsmath}
\usepackage{array}

\definecolor{Navy}{HTML}{344354}
\definecolor{Slate}{HTML}{5A6B7F}
\definecolor{Accent}{HTML}{A85B36}
\definecolor{Pale}{HTML}{EDF1F5}
\definecolor{Ink}{HTML}{242A31}
\setbeamercolor{normal text}{fg=Ink,bg=white}
\setbeamercolor{frametitle}{fg=Navy}
\setbeamercolor{structure}{fg=Navy}
\setbeamercolor{block title}{fg=white,bg=Navy}
\setbeamercolor{block body}{fg=Ink,bg=Pale}
\setbeamercolor{alerted text}{fg=Accent}
\setbeamerfont{frametitle}{size=\Large,series=\bfseries}
\setbeamertemplate{navigation symbols}{}
\setbeamertemplate{itemize item}{\small\raise0.3ex\hbox{\textbullet}}
\setbeamertemplate{footline}{%
  \leavevmode\hbox{\begin{beamercolorbox}[wd=\paperwidth,ht=2.8ex,dp=1.3ex,leftskip=0.5cm,rightskip=0.5cm]{normal text}%
    \scriptsize\color{Slate}Anthropic IPO valuation \hfill 27 Nov 2026 \hfill \insertframenumber/\inserttotalframenumber
  \end{beamercolorbox}}}
\setbeamersize{text margin left=0.6cm,text margin right=0.6cm}
\hypersetup{colorlinks=true,urlcolor=Slate}
\newcommand{\source}[1]{\par\vspace{0.08cm}{\scriptsize\color{Slate}#1}}
\newcommand{\bn}{\,\text{bn}}
\title{Anthropic IPO Valuation}
\subtitle{A basic revenue-multiple model and a DCF cross-check}
\date{27 November 2026}

\begin{document}

% 1 --------------------------------------------------------------------------
\begin{frame}{Anthropic IPO Valuation}
  {\large A basic model: cash flows, growth and the discount rate}
  \vspace{0.35cm}
  \begin{columns}[T,onlytextwidth]
    \begin{column}{0.48\textwidth}
      \begin{block}{Revenue-multiple enterprise value}
        \centering\vspace{0.12cm}
        {\Huge\bfseries\$1.95 trillion}\\[0.15cm]
        \$195bn of 2028 revenue $\times$ 10.0x
        \vspace{0.12cm}
      \end{block}
    \end{column}
    \begin{column}{0.48\textwidth}
      \begin{block}{DCF enterprise value}
        \centering\vspace{0.12cm}
        {\Huge\bfseries\$788.5 billion}\\[0.15cm]
        Ten-year forecast; 12\% WACC; 3\% terminal growth
        \vspace{0.12cm}
      \end{block}
    \end{column}
  \end{columns}
  \vspace{0.18cm}
  \begin{itemize}
    \item Historical investor discussion is about \$1.5tn; the earlier \$2tn is a historical stress case.
    \item The assumed multiple gives a value 30.0\% above that discussion; the DCF is 47.4\% below it.
    \item December net cash remains \textbf{assumed zero}; September financing sensitivities are separate.
  \end{itemize}
  \source{Information cutoff: 27 Nov 2026. Assumed valuation date: 31 Dec 2026. All values in USD. [S2, S25]}
\end{frame}

% 2 --------------------------------------------------------------------------
\begin{frame}{Cash-flow period: 2027--2036, then a terminal value}
  \begin{columns}[T,onlytextwidth]
    \begin{column}{0.54\textwidth}
      \begin{block}{Explicit forecast: ten annual periods}
        \begin{itemize}
          \item Forecast \textbf{unlevered free cash flow} for 2027--2036.
          \item Allow revenue growth and EBIT margins to converge toward maturity.
          \item Start the perpetual-growth phase in \textbf{2037}.
        \end{itemize}
      \end{block}
      \vspace{0.15cm}
      A five-year forecast would end while assumed growth is still 20\%, creating a sharp maturity transition.
    \end{column}
    \begin{column}{0.42\textwidth}
      \begin{tabular}{@{}lr@{}}
        \toprule
        DCF component & PV (\$bn)\\
        \midrule
        2027--2036 cash flows & 337.2\\
        Terminal value & 451.4\\
        \midrule
        \textbf{Enterprise value} & \textbf{788.5}\\
        \bottomrule
      \end{tabular}
      \vspace{0.35cm}
      \textbf{57.2\%} of DCF value comes from cash flows after 2036.
      \vspace{0.25cm}

      The model discounts at year-end: 2027 is $t=1$ and 2036 is $t=10$.
    \end{column}
  \end{columns}
  \source{2026 is a revenue anchor, not a discounted cash-flow period: the assumed valuation date is 31 Dec 2026.}
\end{frame}

% 3 --------------------------------------------------------------------------
\begin{frame}{Annual growth: rapid expansion, followed by a gradual fade}
  \begin{columns}[T,onlytextwidth]
    \begin{column}{0.52\textwidth}
      \small
      \begin{tabular}{@{}lrr@{}}
        \toprule
        Year & Revenue (\$bn) & Growth\\
        \midrule
        2026E & 64.0 & Anchor\\
        2027E & 128.0 & 100.0\%\\
        2028E & \textbf{195.0} & \textbf{52.3\%}\\
        2029E & 263.3 & 35.0\%\\
        2030E & 329.1 & 25.0\%\\
        2031E & 394.9 & 20.0\%\\
        2032E & 454.1 & 15.0\%\\
        2033E & 499.5 & 10.0\%\\
        2034E & 534.5 & 7.0\%\\
        2035E & 561.2 & 5.0\%\\
        2036E & 578.0 & 3.0\%\\
        \bottomrule
      \end{tabular}
    \end{column}
    \begin{column}{0.44\textwidth}
      \begin{block}{How growth is calculated}
        \[
          R_t=R_{t-1}(1+g_t)
        \]
        \[
          g_{2028}=\frac{195}{128}-1=52.3\%
        \]
      \end{block}
      \small
      \textbf{2028 anchor:} midpoint of the reported \$190--200bn company projection. [S2]
      \vspace{0.2cm}

      \textbf{Other years:} analyst assumptions. The \$64bn 2026 estimate is not reported annual revenue.
      \vspace{0.2cm}

      \textbf{Run rate $\ne$ annual revenue.} A recent revenue pace cannot be substituted for full-year recognized sales.
    \end{column}
  \end{columns}
\end{frame}

% 4 --------------------------------------------------------------------------
\begin{frame}{Cash flow: convert revenue into cash available to capital providers}
  \begin{block}{Unlevered free cash flow (UFCF)}
    \centering
    $\mathrm{UFCF}_t=\mathrm{EBIT}_t(1-T)+\mathrm{D\&A}_t-\mathrm{Capex}_t-\Delta\mathrm{NWC}_t$
  \end{block}
  \begin{columns}[T,onlytextwidth]
    \begin{column}{0.48\textwidth}
      \small
      \begin{tabular}{@{}lr@{}}
        \toprule
        Driver & Assumption\\
        \midrule
        EBIT margin & 10\% $\rightarrow$ 35\%\\
        Tax on positive EBIT & 25\%\\
        D\&A / revenue & 3\%\\
        Cash capex / revenue & 8\%\\
        Operating NWC / revenue & 2\%\\
        \bottomrule
      \end{tabular}
      \vspace{0.2cm}

      Operating compute and distribution costs enter assumed EBIT margins; contract timing is not fully modeled.
    \end{column}
    \begin{column}{0.47\textwidth}
      \small
      \begin{tabular}{@{}lrrr@{}}
        \toprule
        \$bn & 2027 & 2028 & 2036\\
        \midrule
        Revenue & 128.0 & 195.0 & 578.0\\
        After-tax EBIT & 9.6 & 29.3 & 151.7\\
        D\&A & 3.8 & 5.9 & 17.3\\
        Capex & (10.2) & (15.6) & (46.2)\\
        $\Delta$NWC & (1.3) & (1.3) & (0.3)\\
        \midrule
        \textbf{UFCF} & \textbf{1.9} & \textbf{18.2} & \textbf{122.5}\\
        \bottomrule
      \end{tabular}
      \vspace{0.2cm}

      Example: 2028 UFCF $=29.25+5.85-15.60-1.34=18.16$.
    \end{column}
  \end{columns}
  \source{Rounded displayed figures. No tax-loss carryforwards or separate funding schedule. Long-term compute commitments are not deducted again as one-year capex.}
\end{frame}

% 5 --------------------------------------------------------------------------
\begin{frame}{Discount rate: 12\% WACC, with a 10--14\% sensitivity}
  \small
  \begin{columns}[T,onlytextwidth]
    \begin{column}{0.52\textwidth}
      \begin{block}{Discounting and terminal value}
        \[
          \mathrm{EV}=\sum_{t=1}^{10}\frac{\mathrm{UFCF}_t}{(1+r)^t}
          +\frac{\mathrm{TV}_{2036}}{(1+r)^{10}}
        \]
        \[
          \mathrm{TV}_{2036}=\frac{\mathrm{UFCF}_{2037}}{r-g}
        \]
      \end{block}
      \small
      \textbf{$r=12\%$:} illustrative WACC for technology, execution and cash-flow risk. It is not an observed or CAPM-estimated company rate.
      \vspace{0.15cm}

      \textbf{$g=3\%$:} mature perpetual revenue growth. Terminal EBIT margin and investment ratios stay at 2036 levels.
    \end{column}
    \begin{column}{0.44\textwidth}
      \small
      \textbf{DCF enterprise value (\$bn)}\\[0.15cm]
      \begin{tabular}{@{}lrrr@{}}
        \toprule
        WACC / $g$ & 2\% & 3\% & 4\%\\
        \midrule
        10\% & 982.1 & 1,074.3 & 1,197.3\\
        12\% & 739.8 & \textbf{788.5} & 849.4\\
        14\% & 581.9 & 610.2 & 644.1\\
        \bottomrule
      \end{tabular}
      \vspace{0.32cm}

      Lower discount rates increase the value of distant cash flows.
      \vspace{0.2cm}

      At zero assumed net cash, even 10\% WACC / 4\% growth remains below the \$1.5tn equity discussion.
    \end{column}
  \end{columns}
  \source{Sensitivity keeps the explicit operating forecast fixed. Terminal revenue, UFCF, investment and discounting recalculate for each rate pair. Require $r>g$.}
\end{frame}

% 6 --------------------------------------------------------------------------
\begin{frame}{What would a \$1.5 trillion discussion value require?}
  \begin{columns}[T,onlytextwidth]
    \begin{column}{0.48\textwidth}
      \begin{block}{Revenue-multiple interpretation}
        \[
          \mathrm{EV}=\$195\bn\times10=\$1{,}950\bn
        \]
        \[
          \frac{\$1{,}500\bn}{\$195\bn}=7.69\mathrm{x}
        \]
      \end{block}
      \small
      At the \$195bn forecast, 8x gives \$1.56tn and 12x gives \$2.34tn.
      \vspace{0.2cm}

      The 10x base multiple is an assumption. It is not a current peer-derived median.
    \end{column}
    \begin{column}{0.48\textwidth}
      \begin{block}{DCF interpretation}
        \small
        Holding the ten-year forecast, WACC and investment ratios fixed:
        \begin{itemize}
          \item Required 2037 UFCF: \textbf{\$325.0bn}.
          \item Required perpetual EBIT margin: \textbf{79.5\%}.
        \end{itemize}
      \end{block}
      \small
      That margin is substantially above the assumed 35\% terminal margin.
      \vspace{0.2cm}

      Supporting \$1.5tn therefore requires a different growth path, earlier cash generation, lower WACC or other material assumption changes.
    \end{column}
  \end{columns}
  \source{Zero December net cash is assumed. September balances are disclosed; conversion and financing sensitivities follow.}
\end{frame}

% 7 --------------------------------------------------------------------------
\begin{frame}{Conclusion and sources}
  \begin{columns}[T,onlytextwidth]
    \begin{column}{0.49\textwidth}
      \begin{block}{Answers to the three model questions}
        \begin{itemize}
          \item \textbf{Cash-flow period:} 2027--2036; terminal value from 2037.
          \item \textbf{Annual growth:} 100\%, 52.3\%, then a fade to 3\% by 2036.
          \item \textbf{Discount rate:} 12\% WACC; examine 10--14\%.
        \end{itemize}
      \end{block}
      \small
      \textbf{Assessment:} \$1.95tn is consistent with the selected revenue multiple. The \$788.5bn DCF shows that cash economics need much stronger assumptions to support the reported discussion value.
    \end{column}
    \begin{column}{0.47\textwidth}
      \footnotesize
      \textbf{Public sources}
      \begin{itemize}
        \item[S1] \href{https://www.anthropic.com/news/series-h}{Anthropic, 28 May 2026}: \$965bn post-money Series H valuation.
        \item[S2] \href{https://www.investing.com/news/stock-market-news/anthropic-ipo-valuation-hinges-on-190200-billion-2028-revenue-forecast-sources-say-4862222?ampMode=1}{Reuters, Aug 2026}: reported \$190--200bn 2028 revenue projection.
        \item[S3] \href{https://www.streetinsider.com/Reuters/Exclusive-Anthropic\%2BIPO\%2Bprospectus\%2Blays\%2Bbare\%2Bdeep\%2Bdependence\%2Bon\%2BBig\%2BTech\%2Bpartners/27123187.html}{Reuters, 29 Sep 2026}: approximately \$2tn target and cloud-partner dependence.
        \item[S25] \href{https://www.sec.gov/Archives/edgar/data/2130357/000162828026067168/anthropicpbc-sx1.htm}{Anthropic public S-1, 26 Oct 2026}: audited annuals, unaudited interim statements and conversion terms.
      \end{itemize}
      \textbf{Limits:} projections may change. September balances are historical; December net cash and future dilution remain uncertain. The final prospectus provides the priced terms.
    \end{column}
  \end{columns}
  \source{Source: accompanying editable Excel model. Forecasts are analyst assumptions.}
\end{frame}

\begin{frame}{Public S-1: operating loss, net loss and cash flow}
  \small
  \begin{tabular}{@{}lrrr@{}}
    \toprule
    US\$bn & FY2025 & H1 2025 & H1 2026\\
    \midrule
    Revenue & 4.585 & 1.163 & 16.286\\
    Operating income (loss) & (8.060) & (2.135) & (0.693)\\
    Net income (loss) & (41.965) & (12.969) & (80.748)\\
    Financing fair-value loss & 34.308 & 10.942 & 72.388\\
    Operating cash flow & (2.637) & (1.513) & 3.929\\
    Cash capex & 0.131 & 0.018 & 0.482\\
    \bottomrule
  \end{tabular}
  \vspace{0.15cm}
  \begin{itemize}
    \item Q1/Q2 revenue is US\$4.731bn / US\$11.555bn. Q3 is US\$16.956bn in the November S-1/A.
    \item The US\$64bn annual assumption requires US\$30.758bn in Q4, or 81.4\% sequential growth.
    \item H1 cash flow less cash capex is US\$3.447bn; it is not the forecast's unlevered FCF.
  \end{itemize}
  \source{S-1 MD\&A p160; financial statements F-7/F-12. FY2025 audited; interim periods unaudited. Non-cash remeasurement and payment timing affect the cash-flow comparison.}
\end{frame}


\begin{frame}{September cash bridge: a historical sensitivity}
  \small
  Liquid resources US\$131.845bn exclude restricted cash. Subtract US\$1.674bn finance leases, then any retained Amazon note principal.
  \vspace{0.2cm}
  \begin{tabular}{@{}lrr@{}}
    \toprule
    US\$bn unless stated & Notes retained & Full conversion\\
    \midrule
    Amazon principal deducted & 6.203 & 0.000\\
    Net cash proxy & 123.968 & 130.171\\
    DCF equity proxy & 912.507 & 918.710\\
    Required terminal EBIT margin & 71.8\% & 71.4\%\\
    \bottomrule
  \end{tabular}
  \vspace{0.2cm}
  \begin{itemize}
    \item Note carrying value is US\$130.502bn; it is not the US\$6.203bn contractual cash principal.
    \item Conversion is subject to Amazon's 12\% ownership cap unless waived; residual notes may remain.
    \item No Q4 cash roll-forward or IPO proceeds. These are not December cash forecasts or values allocated to existing common shares.
  \end{itemize}
  \source{S-1/A MD\&A p168; balance sheet F-6; lease note 7. Unchanged DCF EV US\$788.5bn; margin calculation uses the US\$1.5tn discussion benchmark.}
\end{frame}

\begin{frame}{Commitments and offering terms still constrain valuation}
  \begin{itemize}
    \item September minimum purchase commitments: US\$313.819bn. Leases not yet commenced: US\$451.2bn. Separate multiyear categories, with stated conditions.
    \item The earlier Reuters US\$518bn observation differs in date and scope. It is not a comparable growth series.
    \item Do not deduct all commitments as current debt or one-year capex. Operating rent stays within forecast margins.
    \item The final prospectus provides the offer price and pro forma share counts. See the debut comparison on the final slide.
  \end{itemize}
  \source{\href{https://www.sec.gov/Archives/edgar/data/2130357/000162828026068676/anthropicpbc-sx1ax1.htm}{S-1/A, 2 Nov 2026}, p174--175. SEC HTML checked; annual forecasts remain assumptions.}
\end{frame}

\begin{frame}{November amendment: nine-month results}
  \begin{itemize}
    \item Nine-month revenue US\$33.242bn; operating loss US\$1.260bn.
    \item Net loss US\$104.039bn includes US\$94.324bn financing fair-value loss.
    \item Operating cash flow US\$6.569bn less cash capex US\$0.925bn is US\$5.644bn; not unlevered FCF.
    \item Aggregate commitments, including short-term and subsequent arrangements, are about US\$895.4bn. Do not add this aggregate to its component categories.
    \item Operating forecasts and zero December net cash remain assumptions; preliminary terms are separately disclosed.
  \end{itemize}
  \source{\href{https://www.sec.gov/Archives/edgar/data/2130357/000162828026068676/anthropicpbc-sx1ax1.htm}{S-1/A, 2 Nov 2026}, p161/174--175; F-8/F-13. Interim figures unaudited.}
\end{frame}


\begin{frame}{November 10 preliminary terms}
  \begin{itemize}
    \item Proposed US\$70--78 range; 825m shares, of which 801.351m are issuer shares and 23.649m are secondary.
    \item At US\$74: issuer gross proceeds US\$59.30bn; seller gross US\$1.75bn does not accrue to the company.
    \item Concurrent private placement: US\$26bn. Combined estimated net proceeds: US\$84.6bn, before separate US\$3.8bn RSU withholding.
    \item Assumed Amazon principal conversion: US\$0.5793bn, not all historical principal.
    \item October run-rate US\$81bn is not full-year revenue. The US\$1.5tn benchmark is historical; no forecast uplift.
  \end{itemize}
  \source{\href{https://www.sec.gov/Archives/edgar/data/2130357/000162828026072484/anthropicpbc-sx1ax2.htm}{S-1/A amendment No. 2, November 10, 2026}, cover, Offering and Use of Proceeds. Historical terms before final pricing.}
\end{frame}

\begin{frame}{Final IPO pricing and Nasdaq debut}
  \begin{itemize}
    \item Anthropic priced at US\$83. Its November 19 debut closed at US\$97.17, up 17.1\%.
    \item In the base deal, issuer IPO gross was US\$79.38bn and selling holders' gross was US\$1.96bn.
    \item Including the US\$26bn placement, estimated net proceeds to the issuer were US\$104.8bn, or US\$117.0bn in the full-option case. Estimated RSU tax withholding is US\$4.3bn.
    \item Base pro forma A plus economic N shares imply a US\$1,202.3bn debut equity proxy. The prospectus assumes US\$0.6127bn of Amazon principal converts.
    \item Settlement was outside the previous November 20, 17:30 Beijing snapshot. See the next slide.
  \end{itemize}
  \source{\href{https://www.sec.gov/Archives/edgar/data/2130357/000162828026074944/anthropicpbc-424b4.htm}{Final prospectus, Form 424B4}, and Nasdaq historical prices.}
\end{frame}

\begin{frame}{Settlement and November 25 public close}
  \begin{itemize}
    \item The November 20 US settlement included the full option and US\$26.0bn private placement.
    \item Combined issuer gross proceeds were US\$117.58bn. Sellers separately received US\$1.96bn gross.
    \item Prospectus net proceeds were estimated at US\$117.0bn, or US\$112.7bn after estimated RSU withholding.
    \item ANT closed at US\$82.79 on November 25. The pro forma economic share count implies a US\$1,036.5bn equity proxy.
    \item September cash adjusted for financing gives a US\$1,025.8bn DCF equity sensitivity. The cash bridge excludes later operating cash changes.
  \end{itemize}
  \source{\href{https://www.sec.gov/Archives/edgar/data/2130357/000162828026075973/ant-20261119.htm}{Form 8-K}, final prospectus and Nasdaq.}
\end{frame}

\end{document}
